% BHU PYQ -- Transcribed & Corrected
% Paper: CHEMN-11 / CHEMN11: Basic Concepts of Chemistry - I
% Exam: B.Sc. (Hons.) Semester I Examination, 2024-25 (NEP)
% Category: Minor Course
% Department: Chemistry

\documentclass[11pt,a4paper]{article}
\usepackage[a4paper,top=2.5cm,bottom=2.5cm,left=2.8cm,right=2.8cm,headheight=14pt]{geometry}
\usepackage{amsmath,amssymb}
\usepackage{fontspec}
\setmainfont{Kohinoor Devanagari}
% \usepackage[T1]{fontenc}
% \usepackage[utf8]{inputenc}
\usepackage{lmodern,microtype,fancyhdr,enumitem,hyperref,lastpage,parskip}

\hypersetup{
    colorlinks=true,
    urlcolor=black,
    linkcolor=black,
    pdftitle={CHEMN11: Basic Concepts of Chemistry I -- BHU Sem I 2024-25 (NEP)}
}

\pagestyle{fancy}
\fancyhf{}
\renewcommand{\headrulewidth}{0.4pt}
\renewcommand{\footrulewidth}{0.4pt}
\fancyhead[L]{\small   \textit{Banaras Hindu University}}
\fancyhead[R]{\small   \textit{CHEMN-11: Basic Concepts of Chemistry - I}}
\fancyfoot[C]{\small Page  \thepage\ of \pageref{LastPage}}


\newlist{parts}{enumerate}{2}
\setlist[parts,1]{label=(\alph*),leftmargin=*,itemsep=4pt,topsep=2pt}
\setlist[parts,2]{label=(\roman*),leftmargin=*,itemsep=2pt,topsep=2pt}

\newcommand{\pts}[1]{\hfill   \textit{[#1]}}


\begin{document}


\begin{center}
  {\large\bfseries BANARAS HINDU UNIVERSITY}\\[2pt]
  {\normalsize B.Sc. (Hons.) --- Semester I Examination, 2024-25 (NEP)}\\[6pt]
  \rule{0.8\linewidth}{0.5pt}\\[6pt]
  {\large\bfseries CHEMISTRY}\\[4pt]
  {\large\bfseries Paper Code: CHEMN-11 : Basic Concepts of Chemistry --- I}\\[4pt]
  {\normalsize Time: 2:30 Hours\hfill Maximum Marks: 70}\\[2pt]
  {\small REF NO. CECONF/N/2024-25/29}
\end{center}

\rule{\linewidth}{0.4pt}

\smallskip



\noindent   \textit{Instructions: Answer   \textbf{five} questions in all including Question No.~1 which is compulsory. All questions carry equal marks (14 marks each).}

\bigskip




\noindent   \textbf{Question 1.} Short Answer Questions (Compulsory):\pts{$2  \times 7 = 14$}

\begin{parts}
  \item
  
\begin{parts}
    \item What is the physical significance of $+$ and $-$ signs in atomic wave function orbitals?\pts{1}
    \item What are the essential conditions for the acceptability of wave functions obtained by solving the Schrödinger equation?\pts{2}
  \end{parts}
  \item What are the major consequences of the inert pair effect in $p$-block elements?\pts{2}
  \item Obtain the mathematical expressions for van der Waals constants $a$ and $b$ in terms of critical constants $P_c, V_c, T_c$, and $R$.\pts{3}
  \item State the First Law of Thermodynamics and write its fundamental limitations.\pts{3}
  \item Prove that the work done during isothermal reversible expansion of an ideal gas is greater than that in adiabatic reversible expansion.\pts{3}
\end{parts}

\medskip\rule{\linewidth}{0.2pt}\medskip




\noindent   \textbf{Question 2.}

\begin{parts}
  \item Calculate the effective nuclear charge ($Z^*$) on the outermost valence electron of $^{29} \text{Cu}$ using Slater's rules.\pts{2}
  \item What are eigenfunctions and eigenvalues? Give two examples for each concept.\pts{2}
  \item What are the solubility trends of alkali metal fluorides and iodides in water? Discuss the reasons based on lattice energy versus hydration energy.\pts{4}
  \item What are the major limitations of the Born-Landé equation for lattice energy calculation?\pts{2}
  \item
  
\begin{parts}
    \item Calculate the standard reduction potential for $ \text{Fe}^{3+}/ \text{Fe}$. Given $E^\circ( \text{Fe}^{3+}/ \text{Fe}^{2+}) = +0.77  \text{ V}$ and $E^\circ( \text{Fe}^{2+}/ \text{Fe}) = -0.47  \text{ V}$.\pts{2}
    \item Will $ \text{Fe}^{2+}$ disproportionate into $ \text{Fe}^{3+}$ and $ \text{Fe}$ in acidic solution? Support your answer with thermodynamic calculation.\pts{2}
  \end{parts}
\end{parts}

\medskip\rule{\linewidth}{0.2pt}\medskip




\noindent   \textbf{Question 3.}

\begin{parts}
  \item Derive the 3D time-independent Schrödinger wave equation assuming the electron behaves as a 3D standing de Broglie wave.\pts{4}
  \item Qualitatively discuss the periodic trends of first ionization energies for the elements of the 2nd and 3rd periods, explaining anomalies at B, O, Al, and S.\pts{4}
  \item Discuss the dependence of lattice energy on ionic size and explain the consequences of ionic size on lattice energy and hydration enthalpy.\pts{3}
  \item Calculate the electronegativity of carbon using Pauling's scale if $\Delta$ (extra bond energy for C--H bond) is $24.3  \text{ kJ mol}^{-1}$. Explain the concept of extra bond energy.\pts{3}
\end{parts}

\medskip\rule{\linewidth}{0.2pt}\medskip




\noindent   \textbf{Question 4.}

\begin{parts}
  \item State the fundamental postulates of the kinetic molecular theory of gases, providing experimental evidence for each postulate.\pts{4}
  \item Explain critical temperature ($T_c$), critical pressure ($P_c$), and critical volume ($V_c$) with suitable examples. How is $V_c$ determined experimentally?\pts{5}
  \item State the Law of Corresponding States and derive the van der Waals reduced equation of state.\pts{5}
\end{parts}

\medskip\rule{\linewidth}{0.2pt}\medskip




\noindent   \textbf{Question 5.}

\begin{parts}
  \item The number of drops of a liquid and water were observed to be 66 and 30 respectively using a stalagmometer. Calculate the surface tension of the liquid if the surface tension of water is $72.75  \text{ dyne/cm}$, and the masses of liquid and water for the same volume are $10.750  \text{ g}$ and $11.250  \text{ g}$ respectively.\pts{5}
  \item What is the Law of Rational Indices? Why is it necessary to convert Weiss indices into Miller indices? Write the steps involved in calculating Miller indices.\pts{4}
  \item Draw neat diagrams and explain the crystal structures of graphite and diamond, detailing hybridization, bonding, and physical properties.\pts{5}
\end{parts}

\medskip\rule{\linewidth}{0.2pt}\medskip




\noindent   \textbf{Question 6.}

\begin{parts}
  \item Using the fundamental thermodynamic relation:
  \[
    \mu_{ \text{JT}} = \frac{1}{C_p} \left[ T \left( \frac{\partial V}{\partial T} \right)_P - V \right]
  \]
  show under what conditions the Joule-Thomson coefficient ($\mu_{ \text{JT}}$) is positive (cooling effect) for a real gas.\pts{5}
  \item Why is cooling of real gases observed during Joule-Thomson expansion? Show that Joule-Thomson expansion is an isenthalpic process ($H =  \text{constant}$).\pts{5}
  \item What is the physical significance of entropy? ``The entropy of the universe is constantly increasing.'' Justify this statement using the Second Law of Thermodynamics.\pts{4}
\end{parts}

\medskip\rule{\linewidth}{0.2pt}\medskip




\noindent   \textbf{Question 7.}

\begin{parts}
  \item Prove the Maxwell thermodynamic relation:
  \[
    \left( \frac{\partial S}{\partial V} \right)_T = \left( \frac{\partial P}{\partial T} \right)_V
  \]\pts{5}
  \item ``The total heat absorbed by a system cannot be completely converted into useful work.'' Explain this statement based on the Carnot cycle efficiency.\pts{5}
  \item Show that $-\Delta G = W_{ \text{net}}$ (useful net work). Under what thermodynamic condition does the change in Gibbs free energy ($\Delta G$) equal the change in Helmholtz work function ($\Delta A$)?\pts{4}
\end{parts}

\vfill
\rule{\linewidth}{0.4pt}

\begin{center}\small
\href{https://bhu-exam.vercel.app/}{Click here to visit our website}\\[4pt]
\href{https://me-aryan.vercel.app/}{Click here to visit our contributor}
\end{center}

\end{document}